\section{Transfer on a binary cube}\label{sec:binary-algebra}
\label{sec:h:binary-algebra}

Both analytic arguments will produce one determinant coordinate at each
vertex of a binary cube. The following algebra transfers a bounded
valuation at every nonzero vertex to the zero vertex. Its essential
feature is that the required dimension depends on a bounded number of
coefficients, even when the family involves arbitrarily many primes.

Put $\Lambda=\Ztwo[[t]]$ and
$O=\widehat{\Lambda_{(2)}}$. Let $G\simeq(\Z/2\Z)^b$ and
$R=\Lambda[G]$. Index the sign homomorphisms $G\to\{\pm1\}$
by $x\in\F_2^b$, with $x=0$ the augmentation. For an element or
matrix over $O[G]$, a subscript $x$ denotes evaluation at this character.
The proofs first control Laurent coefficients, then clear determinant
denominators, and finally use the parity of the number of solutions of
low-degree equations on the cube.

\begin{lemma}[Laurent expansions]\label{lem:h:laurent-expansions}
For $M\ge1$,
\[
 O/2^MO\simeq(\Z/2^M\Z)((t)).
\]
Thus an element of $O$ has a unique expansion
$\sum_{n\in\Z}a_nt^n$ with $a_n\in\Ztwo$ and
$v(a_n)\to\infty$ as $n\to-\infty$. Its $O$-valuation is
$\min_n v(a_n)$.
\end{lemma}

\begin{proof}
If $s\in\Lambda\setminus(2)$, write
$s=t^h v_1+2w_1$, with $v_1\in\Lambda^\times$ and
$w_1\in\Lambda$. Modulo $2^M$ its inverse is a geometric expansion
of length $M$ in $2w_1/(t^h v_1)$. It is a Laurent series with finite
negative support. Conversely every Laurent series with finite negative
support lies in the localization of $\Lambda$ modulo $2^M$, since
$t$ is invertible there. Taking inverse limits proves the expansion
statement. Reduction modulo successive powers of $2$ proves the
valuation formula.
\end{proof}

\begin{lemma}[A uniform pole bound for Schur elimination]
\label{lem:h:bounded-poles}
Let $A_1,\ldots,A_m$ be square matrices over $\Lambda$ with nonzero
residual determinants and
$\ord_t(\overline{\det A_j})\le h$. Put
$A=\diag(A_1,\ldots,A_m)$. If $B=A+2E_1$ with $E_1$ a matrix over
$\Lambda$, then $B$ is invertible over $O$. Modulo $2^M$, every entry
of $B^{-1}$ has no powers below $-Mh$, independently of $m$ and of
the matrix sizes. The same bound holds for a Schur complement
$D=P-QB^{-1}S$ when $P,Q,S$ have entries in $\Lambda$.

In particular, suppose a block over $R$ augments to $A$. It is
invertible over $O[G]$, and the assertion applies uniformly at all sign
evaluations. If the lower off-diagonal block $S$ augments to zero, the
Schur complement augments to the old block $P$.
\end{lemma}

\begin{proof}
Write
$\det A_j=t^{h_j}v_j+2w_j$, where $h_j\le h$,
$v_j\in\Lambda^\times$, and $w_j\in\Lambda$. Put
\[
 C_j=\frac{\operatorname{adj}(A_j)}{t^{h_j}v_j},
 \qquad C=\diag(C_1,\ldots,C_m).
\]
Every entry of $C$ has pole at most $h$. We have $BC=1+2F_1$, with
every entry of $F_1$ having pole at most $h$. Consequently
\[
 B^{-1}\equiv C\sum_{j=0}^{M-1}(-2F_1)^j\pmod{2^M}.
\]
The $j$th summand has pole at most $(j+1)h\le Mh$. The number of
summands in a matrix product does not increase a valuation or pole
bound. Multiplication by power-series matrices likewise introduces no
negative powers. This proves the first assertion.

The ring $O[G]$ is local with maximal ideal generated by $2$ and the
augmentation ideal of $G$. The residual determinant of the proposed
block is that of $A$, a unit over $\F_2((t))$, so the block is
invertible. Each sign evaluation differs from augmentation by a
matrix divisible by $2$ with entries in $\Lambda$. Apply the first
assertion to it. The final assertion is the Schur formula with
$S_0=0$.
\end{proof}

\begin{lemma}[A bounded clearing factor]\label{lem:h:pivot-clearing}
Let $D$ be a based bounded complex of finite free $O[G]$-modules, of
total graded rank at most $r$, whose augmentation $D_0$ is a complex
over $\Lambda$. Suppose that the rational specialization of $D_0$ at
$t=0$ has one of the following cohomology patterns:
\begin{enumerate}[label=\textup{(\roman*)}]
 \item dimension one in degrees one and two, and zero elsewhere;
 \item dimension two in degree one, and zero elsewhere.
\end{enumerate}
Suppose that the total torsion length in the integral specialized
cohomology is at most $h_1$.

In case \textup{(i)}, take an integral degree-one cycle $z$ and a closed
functional $\lambda\in2^{-a_0}(D^2)^*$, and suppose each character
constituent over the fraction field of $O$ is either acyclic or has
pattern \textup{(i)}. Let $u_x$ be the determinant tensor of their
cohomology classes, with value zero in the acyclic case. In case
\textup{(ii)}, take two cycles whose denominator exponents sum to at
most $a_0$, suppose each constituent has pattern \textup{(ii)}, and
let $u_x$ be their wedge determinant coordinate. One may multiply these
coordinates by any one element $\eta\in O[G]$.

Then there are $f,d\in O[G]$ with $d_x=f_xu_x$ at every character,
such that $f$ is $2^{a_0}$ times a polynomial of bounded degree in
the differential entries of $D$, the degree depending only on $r$.
Moreover
\[
 f_0\in\Lambda,\qquad f_0(0)\ne0,\qquad
 v(f_0(0))\le C(r,h_1,a_0).
\]
If, at every fixed precision, the entries of $D_x$ have a uniform bound
on negative Laurent support, then so do the $f_x$. No such assertion is
required for $d_x$, for the cycles, or for $\eta$.
\end{lemma}

\begin{proof}
Over $\Ztwo$, put the specialized complex into a sum of free
cohomology terms and one-dimensional two-term complexes with
differentials $2^{a_j}$, by elementary-divisor decomposition. The sum
of the positive $a_j$ is at most $h_1$. Lift the constant integral
basis changes to $D$. Cancel the same pivot positions formally, retaining
their pivot numerators as denominators. The number of operations is bounded
by $r$. Each pivot and basis change is rational in the differential
entries. Let $P_1$ be a product of sufficiently high powers of the
successive pivot numerators so that every denominator in these changes
and in the determinant comparison divides $P_1$. Its degree is bounded
in terms of $r$; its specialization $P_{1,0}(0)$ is nonzero with
valuation bounded in terms of $r,h_1$. Increase the powers once more
if needed to give all final expressions the same denominator.

Over $O[G][1/P_1]$, case \textup{(i)} leaves a complex
$[O[G]\xrightarrow{\delta}O[G]]$ in degrees one and two. Let
$\alpha,\beta$ be the projected cycle and closed functional. Their
determinant tensor, including the cancelled-pivot determinant factor,
is a rational expression in $D,z,2^{a_0}\lambda$. If a field
constituent has $\delta\ne0$, the cycle identity
$\delta\alpha=0$ gives $\alpha=0$, so the same expression yields
the stipulated zero coordinate. In case \textup{(ii)} the remaining
complex is $O[G]^2[-1]$ and its coordinate is the determinant of the
two projected cycles times the pivot determinant factor. In either
case, after denominator clearing, for every constituent on which
$P_{1,x}\ne0$ we have
\[
 u_x=2^{-a_0}\eta_x Q_{1,x}/P_{1,x},\qquad Q_1\in O[G].
\]
Set $f=2^{a_0}P_1^2$ and $d=\eta P_1Q_1$. This gives $d_x=f_xu_x$
on those constituents. Where a pivot numerator first vanishes,
$P_{1,x}=0$, and both sides of the cleared identity are zero. Thus the
identity holds on every character. Since $f$ involves only the
bounded number of differential entries in the pivot calculation,
its degree, augmentation valuation and pole bounds are as asserted.
The factor $\eta$ has been placed in $d$, never in $f$.
\end{proof}

The preceding lemmas isolate the denominator problem: after eliminating
new blocks, the clearing factor uses a complex of bounded size.
The remaining task is to match finitely many coefficients at zero and
at some nonzero vertex. The next lemma gives that match once the cube
has sufficiently many coordinates.

\begin{lemma}[Finite coefficient congruences]
\label{lem:h:boolean-coefficients}
Let $A_1(x),\ldots,A_q(x)\in\Ztwo$ each be one specified Laurent
coefficient of the evaluation of an element of $O[G]$. If
\[
 b>q(2^{M-1}-1),
\]
there is $x\ne0$ such that
$A_i(x)\equiv A_i(0)\pmod{2^M}$ for every $i$.
\end{lemma}

\begin{proof}
For one coefficient write
\[
 A(x)=\sum_{g\in G}a_g(-1)^{g\cdot x}
     =\sum_{S\subseteq\{1,\ldots,b\}}c_Sx_S,
 \qquad x_S=\prod_{i\in S}x_i,
 \qquad 2^{|S|}\mid c_S.
\]
The last divisibility follows by expanding each sign as a product of
$1-2x_i$. The $j$th binary digit of $z\in\Ztwo$ is
$\delta_j(z)=\binom z{2^j}\bmod2$; this holds on nonnegative integers
by binary expansion and then on $\Ztwo$ by continuity. In the ring
of functions on the binary cube with coefficients in $\F_2[[Z]]$,
\[
 (1+Z)^{A(x)}=(1+Z)^{c_\varnothing}
  \prod_{S\ne\varnothing}
    \left(1+x_S\big((1+Z)^{c_S}-1\big)\right).
\]
The nonconstant factor indexed by $S$ starts at degree at least
$2^{|S|}$ in $Z$. A term contributing to $Z^{2^j}$, for $j\ge1$,
therefore has Boolean degree at most
\[
 \sum_S|S|\le\sum_S2^{|S|-1}\le2^{j-1}.
\]
The zeroth digit is constant. The congruences in the statement are
thus equations over $\F_2$ of total degree at most
$q\sum_{j=1}^{M-1}2^{j-1}=q(2^{M-1}-1)$, all vanishing at zero.
For completeness, the number of their common zeros is even: modulo
two it is the sum over the cube of the product of $1+Q$ over all
defining equations $Q$. Every monomial has degree below $b$, hence
omits a variable and sums to zero. This is the
Chevalley--Warning argument \citep{Chevalley1935,Warning1935}. The zero solution therefore has a
nonzero companion.
\end{proof}

We now apply these coefficient congruences to the cleared identity
$d_x=f_xu_x$. Only the bounded negative support of $f_x$ is prescribed;
for $d_x$ we will match its constant coefficient alone. This distinction
keeps the required number of congruences independent of the numerator's
other coefficients.

\begin{lemma}[Transfer through a bounded clearing factor]
\label{lem:h:boolean-transfer}
Let $f,d\in O[G]$ and let $u_x\in\Lambda[1/2]$ satisfy
$d_x=f_xu_x$ at each character. Suppose
$f_0\in\Lambda$ and $c_0=v(f_0(0))<\infty$. Fix $M>c_0$, and
suppose every $f_x\bmod2^M$ has no powers below $-L(M)$. If
\begin{equation}\label{eq:h:cube-threshold}
 b>(L(M)+2)(2^{M-1}-1),
\end{equation}
some $x\ne0$ satisfies
\begin{equation}\label{eq:h:central-transfer}
 v\bigl(u_x(0)-u_0(0)\bigr)\ge M-2c_0.
\end{equation}
Consequently, when $c_0$ and each $L(M)$ have bounds independent of
the family, a uniform upper bound $v(u_x(0))\le B_1$ at all nonzero
vertices forces $v(u_0(0))\le B_1$ for all sufficiently large $b$.
The base value itself is allowed to vary with the family.
\end{lemma}

\begin{proof}
Apply Lemma~\ref{lem:h:boolean-coefficients} to the coefficients of
$f_x$ in degrees $-L(M),\ldots,0$ and to the constant coefficient
of $d_x$. At the supplied nonzero solution, every negative coefficient
of $f_x$ is divisible by $2^M$, and its constant coefficient is
congruent to $f_0(0)$. Therefore its constant coefficient has valuation
$c_0$, and $v_O(f_x)\le c_0$. Since $d_x\in O$, the identity
$d_x=f_xu_x$ gives $v_O(u_x)\ge-c_0$; the same bound holds for
$u_0$. By Lemma~\ref{lem:h:laurent-expansions}, all their coefficients
are in $2^{-c_0}\Ztwo$. Both have zero negative coefficients because
they lie in $\Lambda[1/2]$.

The convergent constant-term convolution gives
\[
 [t^0]d_x\equiv([t^0]f_x)u_x(0)\pmod{2^{M-c_0}},
 \qquad [t^0]d_0=f_0(0)u_0(0).
\]
Combine these identities with the imposed coefficient congruences and
divide by $f_0(0)$. This proves
Equation~\eqref{eq:h:central-transfer}. For the consequence, choose
one integer $M>\max\{c_*,B_1+2c_*\}$, where $c_*$ bounds $c_0$,
and then take $b$ above Equation~\eqref{eq:h:cube-threshold}. If the
base valuation exceeded $B_1$, the transferred nonzero valuation would
also exceed $B_1$, a contradiction. The pivot valuation affects the
required precision and dimension, but not the transferred upper bound.
\end{proof}

\subsection{Finite cochain models}
\label{sec:common-finite-models}

The arithmetic applications also need finite free complexes that retain
maps to the original cochains. The following algebraic construction
provides these models and transports their comparison diagrams.
It applies before any determinant coordinate is chosen.

\begin{lemma}[Finite Artin models]\label{rl:finite-model}
\label{lem:h:finite-models}
Let $(A,\mathfrak m,k)$ be an Artin local ring, and let $D$ be a complex of
flat $A$-modules.  Suppose $H^*(D\otimes_A k)$ is finite dimensional and
supported in $[a,b]$.  There is a bounded finite free complex $C$, with
\[
 \operatorname{rank}_A C^j=\dim_k H^j(D\otimes_A k),
\]
and a strong deformation retraction
\[
 i:C\longrightarrow D,\qquad p:D\longrightarrow C,\qquad
 pi=1,\qquad 1-ip=dh+hd.
\]
Thus maps to the original complex, as well as maps between finite models,
are available.  Any fixed finite collection of chain maps, homotopies,
triangles, and perfect pairings can be transported to these models.
\end{lemma}

\begin{proof}
Flat modules over an Artin local ring are free, including modules of infinite
rank.  Here is a useful version of the argument.  Lift a basis of $M/\mathfrak
mM$ to a flat module $M$, and map the corresponding free module to $M$.
Its cokernel is zero because $\mathfrak m$ is nilpotent.  Flatness shows that
the kernel has zero reduction, and nilpotence again makes the kernel zero.

Split the residual complex as its cohomology, with zero differential, plus
contractible two-term disks.  Lift all the graded bases.  These lifts are
graded isomorphisms: an inverse modulo $\mathfrak m$ lifts to an inverse by
a finite geometric series.  Let $d_0$ be the lifted disk differential and
write $d=d_0+\epsilon$.  All coefficients of $\epsilon$ belong to
$\mathfrak m$.  For the disk contraction $(i_0,p_0,h_0)$, arrange
$h_0^2=h_0i_0=p_0h_0=0$ and
$1-i_0p_0=d_0h_0+h_0d_0$.  The operators
$1+h_0\epsilon$ and $1+\epsilon h_0$ have finite geometric inverses.
The homological perturbation lemma gives explicit contraction maps;
we record the formulas to retain their compatibility with the
arithmetic diagrams \citep[Sections~2.2--2.4 and~3.2]{Crainic2004}.
The perturbation identities give the differential
\[
 d_C=p_0\epsilon(1+h_0\epsilon)^{-1}i_0
\]
on the lifted residual cohomology, and give
\[
 i=(1+h_0\epsilon)^{-1}i_0,\quad
 p=p_0(1+\epsilon h_0)^{-1},\quad
 h=h_0(1+\epsilon h_0)^{-1}.
\]
Substitution, using $(d_0+\epsilon)^2=0$, proves the asserted contraction
identities.  These calculations are valid for infinite disk sums: each
operator is a module homomorphism and every geometric series terminates
at the nilpotence exponent of $\mathfrak m$.

Transport a map $f$ as $p'fi$, retaining the contraction homotopies
for compositions. A quasi-isomorphism between bounded free complexes
has a contractible cone and a homotopy inverse. This transports finite
diagrams, including triangles and perfect dualities.
\end{proof}
